\section{Discussion}
\label{sec:discussion}

\begin{table}[t]
    \centering
    \small
    \renewcommand{\arraystretch}{1.2}
    \begin{tabular}{@{}p{4.3cm}p{2.3cm}p{6.3cm}@{}}
        \toprule
        \textbf{Sub-hypothesis} & \textbf{Verdict} & \textbf{Evidence} \\
        \midrule
        H1a: $r \in [0.4, 0.6]$ on factual and math & Mostly supported & \triviaqa 0.55, \mmlu 0.45, \mathfive 0.61 in range; \gsm 0.37 (few errors limit variance; AUROC 0.85). \\
        H1b: $r < 0.2$ on ethics & Partly supported & \ethjus $r=0.20$ $[0.11, 0.29]$ on the boundary; \ethcs $r=0.33$. Significantly below factual and math, but not near zero. \\
        H2a: vote beats best by 3--7\pp on factual QA & Not supported & $-2.4$ and $-1.8$\pp (n.s.); gains of about 2\pp only on math. \\
        H2b: vote below best on ethics & Supported & $-4.8$\pp (Holm $p=0.008$) and $-4.2$\pp; the deficit grows with ensemble size. \\
        H3: AUROC 0.70--0.80 on factual tasks & Supported (slightly exceeded) & \mmlu 0.78, \triviaqa 0.85, pooled 0.79 $[0.76, 0.82]$. \\
        \bottomrule
    \end{tabular}
    \caption{Verdict on each sub-hypothesis.}
    \label{tab:verdicts}
\end{table}

\Tabref{tab:verdicts} summarizes the verdict on each sub-hypothesis.
We now discuss what the results mean.

\para{Agreement is most informative when the answer space is large.}
When answers are open-ended, as in \triviaqa and \mathfive, independent models rarely produce the same wrong string or number by chance.
Agreement on such tasks is strong evidence of a correct answer, and unanimous answers are right 96--100\% of the time.
Binary answer spaces remove much of this evidence.
Models often agree by chance, and agreement can take only a few values.
Comparisons of agreement across tasks should therefore account for the size of the answer space.
Chance-corrected agreement measures are a natural next step.

\para{Ethics is weaker for two distinct reasons.}
The first reason is format.
The binary factual control loses 0.056 AUROC relative to open and multiple-choice factual tasks, and it matches \ethcs.
The second reason is specific to the domain.
Errors are more correlated on ethics than on the control ($\phi$ 0.45--0.49 vs.\ 0.30), \ethjus is significantly below the control, and consensus errors show shared leanings, such as treating any punishment as wrong.
Our hypothesis attributed the whole ethics deficit to shared bias.
The data suggest that this explains only about half of the gap for commonsense morality, but a larger part for justice.
Part of the measured ``error'' on ethics is also disagreement with crowd labels that are themselves debatable.

\para{Ranking by agreement and voting are different uses of consensus.}
Agreement ranks items well even where voting does not improve accuracy.
On factual QA, knowledge is spread unevenly across models.
The strongest model knows many facts that weaker models do not, and the plurality outvotes it.
H2a implicitly assumed members of similar strength with independent errors.
Our pool has neither property.
In this regime, adding members dilutes the best model on factual tasks, and it reinforces shared mistakes on ethics.
Math is the exception, because wrong chains of thought rarely land on the same answer.

\para{Practical guidance.}
Cross-model agreement is a good label-free {\em triage} signal.
Practitioners can accept unanimous or near-unanimous answers on factual and math tasks, and send low-agreement items to a stronger model or to a human.
It does not replace using the best available model, because a plurality vote over a mixed pool is not a free accuracy gain.
For moral judgements, consensus among models should not stand in for human labels.
Between 10\% and 17\% of unanimous \ethics verdicts contradict the crowd label, and adding models makes the vote worse relative to its best member.
This matters for LLM-as-judge panels and synthetic preference data, where model consensus is often treated as ground truth.

\subsection{Limitations}
\label{sec:limitations}

\para{Model scale and precision.}
All models have 3.8--9B parameters and run at 4-bit precision, because only local open models on shared GPUs were available.
Frontier models may be more accurate, more correlated with each other~\citep{goel2025greatmodels}, or both.
The absolute vote-versus-best numbers in particular may not transfer.

\para{Pool composition.}
Each task has one clearly stronger model, and the vote-versus-best comparison depends strongly on this mix.
The ensemble-size sweep and the one-per-family analysis address this only in part.

\para{Ethics ground truth.}
\ethics labels come from crowdworkers, and the hard test split is adversarially filtered.
Some ``errors'' are label disagreements rather than model mistakes (\secref{sec:res_errors}).
Correctness on ethics therefore partly means agreement with the norms of US crowdworkers.

\para{Post-hoc choices.}
We tightened the \triviaqa equivalence rule after inspecting errors, and we added the binary control after seeing the main results.
Answer matching for \triviaqa and \mathfive is heuristic, and spot checks reveal remaining misses (\eg ``Land Rover'' vs.\ ``Jaguar Land Rover'').
Most such misses lower the measured accuracy of high-agreement items, which makes our factual AUROC estimates conservative.

\para{Prompting and determinism.}
Each model gives one greedy answer, with direct answers on factual and ethics tasks.
Chain of thought on ethics or \mmlu might change the agreement structure, and we did not test prompt variants.
Long 4-bit chains of thought also change with batch composition, so individual \mathfive answers are not reproducible, even though accuracy is stable to within 2\pp.

\para{Contamination and sample size.}
All benchmarks are public and may appear in training data.
This affects absolute accuracy more than the relation between agreement and correctness.
With 500 items per task, confidence intervals on $r$ are about $\pm 0.07$, and vote-versus-best differences of 2\pp cannot be resolved.
